% Part: second-order-logic
% Chapter: syntax-and-semantics
% Section: language-of-sol

\documentclass[../../../include/open-logic-section]{subfiles}

\begin{document}

\olfileid[tr]{sol}{syn}{lan}

\olsection{İkinci Dereceden Mantığın Dili}

Birinci dereceden mantıkta olduğu gibi, ikinci dereceden mantığın ifadeleri
\emph{!!{variable}s}, \emph{!!{constant}s}, \emph{!!{predicate}s} ve bazen
\emph{!!{function}s} içeren temel bir söz dağarcığından kurulur. Bunlardan,
mantıksal bağlaçlar, niceleyiciler ve parantezlerle virgüller gibi noktalama
işaretleriyle birlikte \emph{terimler} ve \emph{!!{formula}s} oluşturulur.
Aradaki fark, ikinci dereceden mantığın nesne değişkenlerine ek olarak bağıntı
ve fonksiyon değişkenleri de içermesi ve bunlar üzerinde nicelemeye izin
vermesidir. Dolayısıyla ikinci dereceden mantığın mantıksal simgeleri, birinci
dereceden mantığın simgelerine ek olarak şunları içerir:

\begin{enumerate}
\item Her $n$ aritesi için, ikinci dereceden bağıntı !!{variable}s{}inden oluşan bir
  !!{denumerable}s küme: $\Obj V_0^n$, $\Obj V_1^n$, $\Obj V_2^n$, \dots
\item Her $n$ aritesi için, ikinci dereceden fonksiyon !!{variable}s{}inden oluşan bir
  !!{denumerable}s küme: $\Obj u_0^n$, $\Obj u_1^n$, $\Obj u_2^n$, \dots
\end{enumerate}

Birinci dereceden mantıkta !!{identity}~$\eq$ genellikle dile katılır.
Birinci dereceden mantıkta bir $\Lang{L}$ dilinin mantıksal olmayan simgeleri,
ilginç herhangi bir şeyi ifade etmemizi sağlamak bakımından çok önemlidir.
Elbette mantıksal olmayan simge kullanmayan !!{sentence}s vardır, fakat
yalnızca~$\eq$ ile ilginç bir şey söylemek zordur. İkinci dereceden mantıkta
ise sınırsız sayıda bağıntı ve fonksiyon değişkenimiz bulunduğundan, özel bir
mantıksal-olmayan simgeler dağarcığı olmadan bile birinci dereceden bir dilde
söyleyebileceğimiz her şeyi söyleyebiliriz.

\end{document}
